% Davian R. Chin — Teaching CV (Mathematics, Computer Science & Physics)
% Compile: tectonic Davian_Chin_CV_Teaching.tex
\documentclass[10pt,a4paper]{article}

\usepackage[a4paper,margin=1.75cm,top=1.5cm,bottom=1.6cm]{geometry}
\usepackage[T1]{fontenc}
\usepackage{lmodern}
\usepackage{microtype}
\usepackage{xcolor}
\usepackage{titlesec}
\usepackage{enumitem}
\usepackage[hidelinks]{hyperref}
\usepackage{ragged2e}

\definecolor{accent}{RGB}{0,96,86}     % teal for the teaching profile
\definecolor{muted}{RGB}{90,90,90}

\pagestyle{empty}
\setlength{\parindent}{0pt}
\setlength{\parskip}{2pt}
\emergencystretch=1.5em

\titleformat{\section}{\large\bfseries\color{accent}}{}{0em}{}[\vspace{2pt}{\color{accent}\titlerule[0.8pt]}]
\titlespacing*{\section}{0pt}{10pt}{7pt}

\newcommand{\entry}[2]{{\textbf{#1}} \hfill {\small\color{muted}#2}\\[1pt]}
\newcommand{\tag}[1]{{\small\color{muted}#1}}
\begin{document}
\RaggedRight

%========================  HEADER  ========================
{\Huge\bfseries\color{accent} Davian R. Chin}\\[3pt]
{\large STEM Educator --- Mathematics, Computer Science \& Physics}\\[5pt]
{\small
\href{mailto:davianc@proton.me}{davianc@proton.me} \,|\, \href{mailto:d.r.chin@pgr.reading.ac.uk}{d.r.chin@pgr.reading.ac.uk} \,|\, +44\,7549\,039\,554 \,|\, London, UK\\
\href{https://www.linkedin.com/in/davianc}{linkedin.com/in/davianc} \,|\,
\href{https://github.com/dave2k77}{github.com/dave2k77} \,|\,
\href{https://orcid.org/0009-0003-9434-3919}{ORCID 0009-0003-9434-3919}}

%========================  PROFILE  ========================
\section{Teaching Profile}
Educator and mathematical scientist with 15+ years' teaching and leadership experience across the UK, China, Kazakhstan and Jamaica (secondary to A level/CAPE). Experienced in curriculum design (CSEC/CAPE/A level), department leadership, and coaching teachers in evidence-informed practice.

\begin{itemize}[leftmargin=1.3em,itemsep=1.5pt,topsep=2pt]
  \item MA Education research: developed and evaluated a computational-thinking framework for mathematics teaching.
  \item Converts advanced STEM ideas into classroom-ready resources: Python/Jupyter lesson libraries, lab-based modules with rubrics, and interactive demonstrators.
  \item Current PhD research (fractional calculus, stochastic processes, scientific machine learning) strengthens subject knowledge across mathematics, physics and computer science.
\end{itemize}

Seeking a teaching position in mathematics, computer science, or physics.


%========================  TEACHING EXPERIENCE  ========================
\section{Teaching \& Leadership Experience}
\entry{Founder \& Director, Neuryte Learning Ltd}{STEM education \& AI consultancy \,--\, 2024--present}
\begin{itemize}[leftmargin=1.3em,itemsep=1.5pt,topsep=2pt]
  \item Consult on the design and delivery of STEM education programmes (mathematics, science, computer science, engineering, AI/ML engineering) for schools, colleges, education providers and edtech organisations.
  \item Lead end-to-end STEM project development and management (scoping, design, delivery, evaluation), including practical AI integration.
  \item Teach and tutor online and face-to-face across mathematics, science, computer science and physics; develop reproducible computational resources in Python (notebooks, demonstrators).
\end{itemize}

\entry{Teacher of Computer Science}{Walton High, Milton Keynes, UK \,--\, Oct 2023--Dec 2024}
\begin{itemize}[leftmargin=1.3em,itemsep=1.5pt,topsep=2pt]
  \item Taught computer science to secondary students preparing for external examinations; built skills in computational thinking, modelling, design, optimisation and data-informed decision-making.
  \item Delivered learning through guided inquiry, team-based real-world projects and hands-on skill practice.
\end{itemize}

\entry{Teacher of Mathematics \& Science}{The Bedford Academy, Bedford, UK \,--\, Sep 2019--Aug 2023}
\begin{itemize}[leftmargin=1.3em,itemsep=1.5pt,topsep=2pt]
  \item Taught mathematics and science through real-world, application-led examples; used technology including simulation and interactive visualisation.
  \item Raised progress for lower-attaining learners through a hands-on, context-based approach; increased STEM engagement through national science competitions.
\end{itemize}

\entry{Mathematics Department Chair \& Lead Teacher}{Shenzhen Vanke Meisha Academy, China \,--\, Aug 2016--Jul 2019}
\begin{itemize}[leftmargin=1.3em,itemsep=1.5pt,topsep=2pt]
  \item Led a multicultural team to build the mathematics department from scratch --- strategic focus, communication protocols and operational framework.
  \item Served on the school's strategic leadership team (curriculum and assessment, school improvement, residential school committees).
  \item Carried out action research on integrating computational thinking and cross-curricular frameworks into mathematics and science teaching.
\end{itemize}

\entry{STEM/CLIL Mathematics Teacher \& Instructional Designer}{Nazarbayev Intellectual Schools, Kazakhstan \,--\, Aug 2014--Jul 2016}
\begin{itemize}[leftmargin=1.3em,itemsep=1.5pt,topsep=2pt]
  \item Taught Cambridge Integrated Mathematics to Grades 11--12 for university entry; supported English-medium learning for Russian- and Kazakh-speaking students and staff.
  \item Delivered masterclasses on effective teaching, learning and assessment; modelled interdisciplinary, problem-based approaches.
  \item Led implementation and evaluation of a summer STEM school on mathematical applications in biodiversity and rocket science.
\end{itemize}

\entry{A-Level Mathematics \& Physics Teacher \& Programmes Coordinator}{St Hugh's High School, Ministry of Education, Jamaica \,--\, Sep 2009--Aug 2014}
\begin{itemize}[leftmargin=1.3em,itemsep=1.5pt,topsep=2pt]
  \item Designed curriculum and instructional plans for CSEC, CAPE and A-Level mathematics and physics courses; taught grades 11--13 in preparation for university entry.
  \item Supervised small mathematics research projects and physics laboratory work; assisted with content and teaching-strategy training for new teachers.
\end{itemize}

\entry{Mathematics Teacher (Grades 10--11)}{Gaynstead High School, Jamaica \,--\, Jan--Aug 2009}

\entry{Curriculum \& lesson design}{Open-source teaching materials \,--\, ongoing}
\begin{itemize}[leftmargin=1.3em,itemsep=1.5pt,topsep=2pt]
  \item \textbf{\href{https://github.com/dave2k77/comp_math}{comp\_math}} --- A-Level mathematics lesson library taught through Python and Jupyter notebooks (rational and nonlinear functions and their graphs, the binomial theorem with tested implementations, statistics and data summarisation).
  \item \textbf{\href{https://github.com/dave2k77/differentiable-probabilistic-programming-curriculum}{differentiable-probabilistic-programming-curriculum}} --- a full curriculum spine for differentiable and probabilistic programming (JAX/NumPyro): modules with prerequisites and learning outcomes, labs, notebooks, capstone projects, assessment rubrics and machine-validated curriculum metadata.
  \item \textbf{Doctoral-era teaching materials} --- guided learning manuals and workbooks on mathematical methods, scientific machine learning and neural networks, signal processing and wavelets, non-Gaussian stable laws, and fractional reservoir computing; plus a scheme of work for a fractional-calculus/long-range-dependence study week.
\end{itemize}

\entry{Educational writing \& public engagement}{2014--present}
\begin{itemize}[leftmargin=1.3em,itemsep=1.5pt,topsep=2pt]
  \item Educational writing and reflection on modern pedagogy, including the essay \emph{A thought on teaching students in the modern world} (2014).
  \item Public engagement with research: conference talks and posters communicating mathematics and critical brain science to non-specialist academic audiences.
\end{itemize}

%========================  SUBJECTS  ========================
\section{Subjects \& Subject Knowledge}
\textbf{Mathematics:} real and complex analysis, calculus, linear algebra, ordinary and partial differential equations, probability and statistics, number theory, discrete and decision mathematics, mathematical modelling.\\
\textbf{Physics:} classical dynamics, heat/diffusion and wave PDEs, statistical physics, dynamical systems and criticality, computational and numerical modelling.\\
\textbf{Computer Science:} programming (Python primary; also C/C++, R, MATLAB, SQL, Haskell), software engineering, databases and networks, artificial intelligence and machine learning, data science, cloud and big-data fundamentals.

%========================  EDUCATIONAL WRITING  ========================
\section{Educational Research \& Writing}
\begin{itemize}[leftmargin=1.3em,itemsep=2pt,topsep=2pt]
  \item MA Education dissertation (University of South Wales, 2018): \emph{Strategies towards a more visible, practical and culturally relevant mathematics education framework: teaching mathematics through computational thinking} --- comparative analysis of formative and summative assessment evidence showing improved outcomes under the computational-thinking approach.
  \item Author of research-adjacent teaching materials: interactive neural-network visualiser for teaching deep learning (\href{https://github.com/dave2k77/neural-net-model}{neural-net-model}, JavaScript), and the lesson libraries and workbooks above.
\end{itemize}

%========================  EDUCATION  ========================
\section{Education}
\entry{PhD, Biomedical Engineering (Computational Neuroscience)}{University of Reading, UK \,--\, 2025--present (expected 2028)}
Research-grade specialisation underpinning subject knowledge: fractional calculus, stochastic processes, functional analysis, machine learning.

\entry{MSc, Computer Science (AI, ML \& Big Data Technologies)}{University of East London, UK \,--\, 2023}
\entry{MA, Education (Development, Management \& Technology)}{University of South Wales, UK \,--\, 2018}
\entry{BSc, Mathematics \& Computer Science}{University of the West Indies, Mona, Jamaica \,--\, 2008}
{\small Graduate coursework and research in mathematics and theoretical physics (UWI Mona, 2010--2013).}

%========================  SKILLS  ========================
\newpage % Keep the technical-skills paragraph together in the downloadable CV.
\section{ICT \& Technical Skills}
\textbf{Teaching technology:} Python, Jupyter notebooks, interactive demonstrators, LaTeX for teaching materials, Google Colab and cloud classrooms, Git/GitHub for versioned coursework, automated curriculum validation (JSON Schema, CI).\\
\textbf{Subject software:} NumPy/SciPy, pandas, scikit-learn, JAX, PyTorch, Matplotlib, MATLAB, R, SQL.\\
\textbf{Systems:} Linux, Docker basics, Google Cloud, high-performance computing.

%========================  CPD  ========================
\section{Certifications \& Professional Development}
\begin{itemize}[leftmargin=1.3em,itemsep=1.5pt,topsep=2pt]
  \item PyTorch Essential Training: Deep Learning (LinkedIn Learning, 2022).
  \item Python Programmer Track (DataCamp, 2020); Python for Data Science: Fundamentals \& Intermediate (Dataquest, 2019).
  \item Computer Science Career Path \& Code Foundations (Codecademy, 2020).
  \item Fellow, Institute of Mathematics and its Applications (FIMA).
\end{itemize}

\vspace{4pt}
{\small\color{muted} References available on request.}

\end{document}
